\section{Introduction and exact scope}
\label{sec:introduction}

The preceding finite-conductor transfer result separates a bounded periodic
second coefficient from the energy-bearing remainder of $\Lambda-1$.  The
periodic part can be passed through $\Lambda(mn+h)$ in a one-sided
Bombieri--Vinogradov range, but low-period spaces capture only a vanishing
fraction of the $L^2$ energy of the prime fluctuation
\citep{WangTPCSeven2026,WangTPCEight2026}.  The next question is therefore
not whether that remainder is large---it is---but whether a specified
prime-target form can nevertheless average it.

This paper answers that question in one nonlocal coordinate.  Fix
$h\in\Z\setminus\{0\}$ and $W\in C_c^1((0,\infty))$.  We extend
$\Lambda(t)$ by zero for $t\leq0$ and define
\begin{equation}\label{eq:intro-complete-residual}
 \cA_h(X;W)
 =\sum_{m,n\geq1}
   (\Lambda(n)-1)\Lambda(mn+h)W(mn/X).
\end{equation}
For $D\geq1$, put
\begin{align}
 \cA_{h,\leq D}(X;W)
 &=\sum_{\substack{m,n\geq1\\n\leq D}}
   (\Lambda(n)-1)\Lambda(mn+h)W(mn/X),
 \label{eq:intro-small-residual}\\
 \cA_{h,>D}(X;W)
 &=\sum_{\substack{m,n\geq1\\n>D}}
   (\Lambda(n)-1)\Lambda(mn+h)W(mn/X).
 \label{eq:intro-large-residual}
\end{align}
The split is exact:
\begin{equation}\label{eq:intro-exact-split}
 \cA_h=\cA_{h,\leq D}+\cA_{h,>D}.
\end{equation}

The form \eqref{eq:intro-complete-residual} does not localize $m$ and $n$
separately.  It fixes their product scale and sums over every factorization.
That distinction is the source of the theorem.  If $r=mn$, then
\begin{equation}\label{eq:intro-zero-mode-collapse}
 \sum_{n\mid r}(\Lambda(n)-1)
 =\log r-\tau(r).
\end{equation}
Thus the full factor coordinate collapses to the difference between an
elementary logarithmic weight and the divisor function.  The target remains
prime-sensitive, but it is now precisely a Titchmarsh divisor problem.

\subsection{The fixed-shift result}

Set
\begin{equation}\label{eq:intro-moments-W}
 I_0(W)=\int_0^\infty W(u)\,du,
 \qquad
 I_1(W)=\int_0^\infty W(u)\log u\,du,
\end{equation}
and define the Euler product
\begin{equation}\label{eq:intro-Chs}
 C_h(s)=
 \prod_{p\mid h}(1-p^{-s-1})
 \prod_{p\nmid h}
 \left(1+\frac1{(p-1)p^{s+1}}\right).
\end{equation}
It is holomorphic in a neighborhood of $s=0$.  We abbreviate
\begin{equation}\label{eq:intro-Ch-notation}
 C_h=C_h(0),\qquad C_h'=C_h'(0).
\end{equation}
Here and below, $\gamma$ denotes the Euler--Mascheroni constant.

Our main theorem, stated precisely in \cref{thm:large-factor-tail}, gives
\begin{align}
 \cA_{h,>D}(X;W)
={}&X\Big[(1-C_h)\{(\log X)I_0(W)+I_1(W)\}
       -2(\gamma C_h+C_h')I_0(W)\Big]
 \notag\\
 &-XI_0(W)
 \sum_{\substack{n\leq D\\(n,h)=1}}
 \frac{\Lambda(n)-1}{\varphi(n)}
 +O_{A,h,W}\!\left(\frac{X}{(\log X)^A}\right),
 \label{eq:intro-tail-formula}
\end{align}
uniformly for
\begin{equation}\label{eq:intro-D-range}
 2\leq D\leq X^{1/2}(\log X)^{-B},
\end{equation}
where $B=B(A,h,W)$ is sufficiently large.  An elementary Dirichlet-series
calculation then shows
\begin{equation}\label{eq:intro-coefficient-sum}
 \sum_{\substack{n\leq D\\(n,h)=1}}
 \frac{\Lambda(n)-1}{\varphi(n)}
 =(1-C_h)\log D+\kappa_h+o_h(1)
\end{equation}
for an explicit finite constant $\kappa_h$.  Hence the leading tail is
\begin{equation}\label{eq:intro-tail-leading}
 (1-C_h)XI_0(W)\log(X/D).
\end{equation}

This is a fixed-shift theorem, including $h=2$.  It evaluates the complete
range beyond the direct Bombieri--Vinogradov cutoff, but only after all of
those factor scales have been added with their original signs.  No estimate
for an individual dyadic $n$-layer follows.

\subsection{What is new and what is imported}

The analytic estimates used here are established results.  The complete
sum is evaluated by the uniform Titchmarsh theorem of
\citet{AssingBlomerLi2021}, whose work supplies strong errors and substantial
uniformity in the shift.  Earlier decisive forms of the Titchmarsh divisor
problem include work of \citet{Fouvry1985} and
\citet{BombieriFriedlanderIwaniec1986}.  The range $n\leq D$ is evaluated by
the classical maximal Bombieri--Vinogradov theorem
\citep{Bombieri1965,Vinogradov1965}.  We prove no improvement to either
input.

The contribution is instead the exact residual coordinate and its boundary:
\begin{enumerate}[label=\textup{(\roman*)}]
 \item the full prime fluctuation $\Lambda(n)-1$, not a bounded periodic
 approximation, is inserted into a genuine fixed-shift prime target;
 \item complete factor aggregation converts it exactly into
 $\log r-\tau(r)$;
 \item the small-factor Bombieri--Vinogradov evaluation is subtracted from
 the complete Titchmarsh evaluation, yielding the entire large-factor tail;
 \item Mellin twisting identifies why only the zero scale frequency
 collapses and why factor localization reintroduces the unresolved bilinear
 geometry.
\end{enumerate}

This distinction parallels a general lesson of the asymptotic sieve:
distribution information and the bilinear information needed to overcome
parity are separate inputs \citep{FriedlanderIwaniec1998}.  Recent work on
prime-producing sieves makes the Type~I/Type~II ledger still more explicit
\citep{FordMaynard2024}.  Our theorem resolves one aggregated coordinate; it
does not verify the bilinear hypotheses required in those frameworks.

For comparison, prime-pair asymptotics are known for almost all shifts in
long windows, with the strongest two-point range relevant here due to
\citet{MatomakiRadziwillTao2019}.  Newer higher-uniformity results handle
arbitrary fixed one-parameter prime-shift patterns $n+a_i h$ in an
almost-all-$h$ regime
\citep{MatomakiRadziwillShaoTaoTeravainen2026}.  Neither statement selects a
prescribed fixed shift such as $h=2$.

\subsection{Strict claim boundary}

The following statements are not consequences of this paper.
\begin{enumerate}[label=\textbf{B\arabic*.},leftmargin=3.1em]
 \item We do not prove an asymptotic for
$\sum_n\Lambda(n)\Lambda(n+h)W(n/X)$ at any prescribed admissible even
shift $h$; in particular, we do not treat $h=2$.
 \item We do not estimate a fixed dyadic box
 $m\asymp M$, $n\asymp N$ with two unrestricted coefficients.
 \item The large-factor theorem permits cancellation between different
 factor scales.  It supplies neither pointwise-in-scale control nor an
 $L^2$ estimate in the scale variable.
 \item No new level of distribution, dispersion theorem, Type~II estimate,
 twin-prime lower bound, or breach of the sieve parity barrier is asserted.
 \item The phrase ``zero Mellin mode'' refers to multiplicative factor-scale
 aggregation.  It is not a spectral realization of the primes or an
 isomorphism with a nonlinear map.
\end{enumerate}

The value of the result is precisely that both sides of the boundary are
rigorous: the zero scale frequency closes at every fixed $h$, while the
localization step that would isolate a prime-pair layer is displayed rather
than assumed.
